\subsubsection{Authentication Controller (authController)}

The Auth Controller is responsible for user authentication, including registration (by a librarian), login, and logout. These tests ensure that access control and user session management are functioning correctly, as specified in the SRS.

\begin{longtable}{|p{0.1\textwidth}|p{0.15\textwidth}|p{0.15\textwidth}|p{0.20\textwidth}|p{0.25\textwidth}|}
    \caption{Test Cases for authController}\\
    \hline
    \textbf{Test Case ID} & \textbf{Related Requirement (SRS)} & \textbf{Test Scenario} & \textbf{Test Steps} & \textbf{Expected Result} \\
    \hline
    \endfirsthead
    \multicolumn{5}{c}%
    {{\bfseries \tablename\ \thetable{} -- continued from previous page}} \\
    \hline
    \textbf{Test Case ID} & \textbf{Related Requirement (SRS)} & \textbf{Test Scenario} & \textbf{Test Steps} & \textbf{Expected Result} \\
    \hline
    \endhead
    \hline \multicolumn{5}{|r|}{{Continued on next page}} \\ \hline
    \endfoot
    \hline
    \endlastfoot

    % Test Cases for registerUser
    TC-AUTH-001 & 3.1.1 User Registration , UC-004  & Successful user registration & 1. Mock \texttt{User.findOne} to resolve as null. \newline 2. Mock \texttt{User.create} to resolve with a new user object. \newline 3. Call \texttt{registerUser} with valid new user data. & 1. Response status is 201. \newline 2. Response JSON indicates successful registration. \\
    \hline
    TC-AUTH-002 & UC-004 (A1: Email Already Exists)  & Attempt to register a user that already exists & 1. Mock \texttt{User.findOne} to resolve with an existing user object. \newline 2. Call \texttt{registerUser} with data of an existing user. & 1. Response status is 400. \newline 2. Response JSON indicates that the username or email already exists. \\
    \hline

    % Test Cases for loginUser
    TC-AUTH-003 & 3.1.2 User Login , UC-001  & Successful user login with valid credentials & 1. Mock \texttt{passport.authenticate} to succeed and return a user object. \newline 2. Call \texttt{loginUser} middleware with correct username and password. & 1. \texttt{req.logIn} is called. \newline 2. Response status is 200. \newline 3. Response JSON indicates successful login and includes user data. \\
    \hline
    TC-AUTH-004 & UC-001 (A1: Invalid Credentials)  & Attempt to login with invalid credentials & 1. Mock \texttt{passport.authenticate} to fail (returns false). \newline 2. Call \texttt{loginUser} middleware with an incorrect username or password. & 1. Response status is 401. \newline 2. Response JSON indicates incorrect credentials. \\
    \hline

    % Test Cases for logoutUser
    TC-AUTH-005 & 3.1.3 User Logout  & Successful user logout & 1. Mock \texttt{req.logout}. \newline 2. Call \texttt{logoutUser}. & 1. \texttt{req.logout} is called. \newline 2. Response status is 200. \newline 3. Response JSON indicates successful logout. \\
    \hline
\end{longtable}

\subsubsection{Author Controller (authorController)}


The Author Controller manages CRUD operations for author entities. These tests validate the creation and retrieval of authors.

\begin{longtable}{|p{0.1\textwidth}|p{0.2\textwidth}|p{0.2\textwidth}|p{0.25\textwidth}|p{0.25\textwidth}|}
    \caption{Test Cases for authorController}\\
    \hline
    \textbf{Test Case ID} & \textbf{Related Requirement (SRS)} & \textbf{Test Scenario} & \textbf{Test Steps} & \textbf{Expected Result} \\
    \hline
    \endfirsthead
    \multicolumn{5}{c}%
    {{\bfseries \tablename\ \thetable{} -- continued from previous page}} \\
    \hline
    \textbf{Test Case ID} & \textbf{Related Requirement (SRS)} & \textbf{Test Scenario} & \textbf{Test Steps} & \textbf{Expected Result} \\
    \hline
    \endhead
    \hline \multicolumn{5}{|r|}{{Continued on next page}} \\ \hline
    \endfoot
    \hline
    \endlastfoot

    % Test Cases for createAuthor
    TC-ATHR-001 & 3.4 Author Management (POST /api/author/add)  & Successfully create a new author & 1. Provide author details in \texttt{req.body}. \newline 2. Mock \texttt{Author.create} to resolve with the new author data. \newline 3. Call \texttt{createAuthor}. & 1. Response status is 201. \newline 2. Response JSON contains the newly created author. \\
    \hline

    % Test Cases for getAllAuthors
    TC-ATHR-002 & 3.4 Author Management (GET /api/author/getAll)  & Successfully retrieve all authors & 1. Mock \texttt{Author.find} to resolve with an array of author objects. \newline 2. Call \texttt{getAllAuthors}. & 1. Response status is 200. \newline 2. Response JSON contains the array of all authors. \\
    \hline
\end{longtable}

\subsubsection{Book Controller (bookController)}

The Book Controller handles all CRUD operations for book entities. The tests ensure that books can be created, retrieved (all and by ID), updated, and deleted as per the requirements.

\begin{longtable}{|p{0.1\textwidth}|p{0.2\textwidth}|p{0.2\textwidth}|p{0.25\textwidth}|p{0.25\textwidth}|}
    \caption{Test Cases for bookController}\\
    \hline
    \textbf{Test Case ID} & \textbf{Related Requirement (SRS)} & \textbf{Test Scenario} & \textbf{Test Steps} & \textbf{Expected Result} \\
    \hline
    \endfirsthead
    \multicolumn{5}{c}%
    {{\bfseries \tablename\ \thetable{} -- continued from previous page}} \\
    \hline
    \textbf{Test Case ID} & \textbf{Related Requirement (SRS)} & \textbf{Test Scenario} & \textbf{Test Steps} & \textbf{Expected Result} \\
    \hline
    \endhead
    \hline \multicolumn{5}{|r|}{{Continued on next page}} \\ \hline
    \endfoot
    \hline
    \endlastfoot

    % Test Cases
    TC-BOOK-001 & 3.3 Book Management (POST /api/book/add) , UC-002  & Successfully create a new book & 1. Provide book details in \texttt{req.body}. \newline 2. Mock \texttt{Book.create} to resolve successfully. \newline 3. Call \texttt{createBook}. & 1. Response status is 201. \newline 2. Response JSON contains the new book. \\
    \hline
    TC-BOOK-002 & UC-002 (A2: API Error)  & Fail to create a book due to a database error & 1. Mock \texttt{Book.create} to reject with an error. \newline 2. Call \texttt{createBook}. & 1. Response status is 500. \newline 2. Response JSON contains an error message. \\
    \hline
    TC-BOOK-003 & 3.3 Book Management (GET /api/book/getAll)  & Successfully retrieve all books & 1. Mock \texttt{Book.find().populate()} to resolve with an array of books. \newline 2. Call \texttt{getAllBooks}. & 1. Response status is 200. \newline 2. Response JSON contains the array of books. \\
    \hline
    TC-BOOK-004 & 3.3 Book Management (GET /api/book/get/:id)  & Successfully retrieve a single book by its ID & 1. Provide a book ID in \texttt{req.params}. \newline 2. Mock \texttt{Book.findById().populate()} to resolve with a book object. \newline 3. Call \texttt{getBookById}. & 1. Response status is 200. \newline 2. Response JSON contains the specific book. \\
    \hline
    TC-BOOK-005 & 3.3 Book Management (GET /api/book/get/:id)  & Attempt to retrieve a book with a non-existent ID & 1. Provide a non-existent book ID in \texttt{req.params}. \newline 2. Mock \texttt{Book.findById().populate()} to resolve as null. \newline 3. Call \texttt{getBookById}. & 1. Response status is 404. \newline 2. Response JSON contains a 'Book not found' message. \\
    \hline
    TC-BOOK-006 & 3.3 Book Management (PUT /api/book/update/:id) & Successfully update an existing book & 1. Provide a book ID and update data. \newline 2. Mock \texttt{Book.findByIdAndUpdate} to resolve with the updated book. \newline 3. Call \texttt{updateBook}. & 1. Response status is 200. \newline 2. Response JSON contains the updated book data. \\
    \hline
    TC-BOOK-007 & 3.3 Book Management (DELETE /api/book/delete/:id)  & Successfully delete a book & 1. Provide a book ID to be deleted. \newline 2. Mock \texttt{Book.findByIdAndDelete} to resolve successfully. \newline 3. Call \texttt{deleteBook}. & 1. Response status is 200. \newline 2. Response JSON contains a 'Book removed' message. \\
    \hline
\end{longtable}

\subsubsection{Borrowal Controller (borrowalController)}

This controller manages borrowal records. Tests cover creating a borrowal (and updating book availability), updating a borrowal's status, and retrieving borrowals based on user roles (Librarian vs. Member).

\begin{longtable}{|p{0.1\textwidth}|p{0.2\textwidth}|p{0.2\textwidth}|p{0.25\textwidth}|p{0.25\textwidth}|}
    \caption{Test Cases for borrowalController}\\
    \hline
    \textbf{Test Case ID} & \textbf{Related Requirement (SRS)} & \textbf{Test Scenario} & \textbf{Test Steps} & \textbf{Expected Result} \\
    \hline
    \endfirsthead
    \multicolumn{5}{c}%
    {{\bfseries \tablename\ \thetable{} -- continued from previous page}} \\
    \hline
    \textbf{Test Case ID} & \textbf{Related Requirement (SRS)} & \textbf{Test Scenario} & \textbf{Test Steps} & \textbf{Expected Result} \\
    \hline
    \endhead
    \hline \multicolumn{5}{|r|}{{Continued on next page}} \\ \hline
    \endfoot
    \hline
    \endlastfoot

    % Test Cases
    TC-BORR-001 & 3.6 Borrowal Management , UC-003  & Create a borrowal for an available book & 1. Mock \texttt{Book.findById} to return an available book. \newline 2. Mock \texttt{Borrowal.create} to succeed. \newline 3. Call \texttt{createBorrowal}. & 1. Book's availability is set to false. \newline 2. Response status is 201. \newline 3. Response JSON contains the new borrowal record. \\
    \hline
    TC-BORR-002 & UC-003 (A1: Book Not Available)  & Attempt to create a borrowal for an unavailable book & 1. Mock \texttt{Book.findById} to return an unavailable book. \newline 2. Call \texttt{createBorrowal}. & 1. Response status is 400. \newline 2. Response JSON has a "Book is not available" message. \\
    \hline
    TC-BORR-003 & 3.6 Borrowal Management (PUT /api/borrowal/update/:id)  & Update borrowal status to 'Returned' & 1. Mock \texttt{Borrowal.findById} to find a borrowal. \newline 2. Mock \texttt{Book.findByIdAndUpdate} to succeed. \newline 3. Call \texttt{updateBorrowal} with status 'Returned'. & 1. Book's availability is set to true. \newline 2. Response status is 200. \newline 3. Response JSON contains the updated borrowal. \\
    \hline
    TC-BORR-004 & 3.6 Borrowal Management (GET /api/borrowal/getAll)  & Get all borrowals as a Librarian & 1. Set \texttt{req.user.role} to 'Librarian'. \newline 2. Mock \texttt{Borrowal.find} to return all borrowals. \newline 3. Call \texttt{getBorrowals}. & 1. \texttt{Borrowal.find} is called with an empty filter. \newline 2. Response contains all borrowal records. \\
    \hline
    TC-BORR-005 & 3.6 Borrowal Management , UC-005  & Get own borrowals as a Member & 1. Set \texttt{req.user.role} to 'Member'. \newline 2. Mock \texttt{Borrowal.find} to return member-specific borrowals. \newline 3. Call \texttt{getBorrowals}. & 1. \texttt{Borrowal.find} is called with a filter for the member's ID. \newline 2. Response contains only the member's borrowals. \\
    \hline
\end{longtable}

\subsubsection{Genre Controller (genreController)}

The Genre Controller manages CRUD operations for genre entities. These tests validate the creation and retrieval of genres.

\begin{longtable}{|p{0.1\textwidth}|p{0.2\textwidth}|p{0.2\textwidth}|p{0.25\textwidth}|p{0.25\textwidth}|}
    \caption{Test Cases for genreController}\\
    \hline
    \textbf{Test Case ID} & \textbf{Related Requirement (SRS)} & \textbf{Test Scenario} & \textbf{Test Steps} & \textbf{Expected Result} \\
    \hline
    \endfirsthead
    \multicolumn{5}{c}%
    {{\bfseries \tablename\ \thetable{} -- continued from previous page}} \\
    \hline
    \textbf{Test Case ID} & \textbf{Related Requirement (SRS)} & \textbf{Test Scenario} & \textbf{Test Steps} & \textbf{Expected Result} \\
    \hline
    \endhead
    \hline \multicolumn{5}{|r|}{{Continued on next page}} \\ \hline
    \endfoot
    \hline
    \endlastfoot

    % Test Cases for createGenre
    TC-GENR-001 & 3.5 Genre Management (POST /api/genre/add)  & Successfully create a new genre & 1. Provide genre details in \texttt{req.body}. \newline 2. Mock \texttt{Genre.create} to resolve with the new genre data. \newline 3. Call \texttt{createGenre}. & 1. Response status is 201. \newline 2. Response JSON contains the newly created genre. \\
    \hline

    % Test Cases for getAllGenres
    TC-GENR-002 & 3.5 Genre Management (GET /api/genre/getAll)  & Successfully retrieve all genres & 1. Mock \texttt{Genre.find} to resolve with an array of genre objects. \newline 2. Call \texttt{getAllGenres}. & 1. Response status is 200. \newline 2. Response JSON contains the array of all genres. \\
    \hline
\end{longtable}

\subsubsection{Review Controller (reviewController)}

The Review Controller is responsible for managing book reviews. The tests cover the creation of reviews and fetching all reviews for a specific book.

\begin{longtable}{|p{0.1\textwidth}|p{0.2\textwidth}|p{0.2\textwidth}|p{0.25\textwidth}|p{0.25\textwidth}|}
    \caption{Test Cases for reviewController}\\
    \hline
    \textbf{Test Case ID} & \textbf{Related Requirement (SRS)} & \textbf{Test Scenario} & \textbf{Test Steps} & \textbf{Expected Result} \\
    \hline
    \endfirsthead
    \multicolumn{5}{c}%
    {{\bfseries \tablename\ \thetable{} -- continued from previous page}} \\
    \hline
    \textbf{Test Case ID} & \textbf{Related Requirement (SRS)} & \textbf{Test Scenario} & \textbf{Test Steps} & \textbf{Expected Result} \\
    \hline
    \endhead
    \hline \multicolumn{5}{|r|}{{Continued on next page}} \\ \hline
    \endfoot
    \hline
    \endlastfoot

    % Test Cases for createReview
    TC-REV-001 & 3.8 Review Management (POST /api/review/add)  & Successfully create a review for a book & 1. Provide review details (bookId, rating, comment) in \texttt{req.body}. \newline 2. Mock \texttt{Review.create} to resolve with the new review object. \newline 3. Call \texttt{createReview}. & 1. Response status is 201. \newline 2. Response JSON contains the newly created review. \\
    \hline

    % Test Cases for getReviewsForBook
    TC-REV-002 & 3.8 Review Management  & Get all reviews for a specific book & 1. Provide a book ID in \texttt{req.params.bookId}. \newline 2. Mock \texttt{Review.find().populate()} to resolve with an array of reviews. \newline 3. Call \texttt{getReviewsForBook}. & 1. \texttt{Review.find} is called with a filter for the book ID. \newline 2. Response status is 200. \newline 3. Response JSON contains the array of reviews for that book. \\
    \hline
\end{longtable}